% ===== 原文 pp.7–8（PDF p.17–18） =====

% [原文 p.7]

\section{主要结果}

在此我们列出本论文给出的主要结果。其中一些结果首先出现在论文 [4] 中，一些出现在论文 [5] 中，
还有一些是本论文的原创工作。
\begin{itemize}[leftmargin=1.8em,itemsep=3pt,topsep=3pt]
\item 在第 3 章，我们定义一个新的四维拓扑规范理论：带规范场的 4d A-模型，其靶空间是
  一个一般黎曼流形 $X$，其上配有李群 $G$ 的作用。
\item 在第 3 章，我们把 Vafa-Witten 扭转 TFT 定位为上述族的特殊成员，即取 $X=\fg$、
  配以 $G$ 在 $\fg$ 上的伴随作用。
\item 在第 5 章，我们把 Vafa-Witten 理论以及 $t=0$、$t=1$、$t=\ii$ 的 GL-扭转理论
  （对非阿贝尔规范群）在圆上紧化。（GL-扭转理论的阿贝尔约化曾在 [4] 中给出。）
\item 在附录 A 与 B 中，我们写下并研究了若干带规范场的 2d 与 3d TFT，并描述了它们的性质。
  特别地，我们证明 2d 带规范场 B-模型的膜范畴是凝聚层的等变导出范畴。
  （这些首先出现在 [4] 的附录中。）
\item 在第 6 章，我们确定 $t=\ii$ 时 GL-扭转理论的线算子的完整幺半范畴
  （它把 [1] 中的 Wilson 线作为特殊情形包含在内）。（首先出现在 [4] 中。）
\item 在第 6 章，我们确定 $t=\ii$ 理论的面算子的完整 2-范畴。（首先出现在 [4] 中。）
\item 在第 6 章，对阿贝尔规范群，我们确定 $t=0$ 与 $t=1$ 理论的面算子 2-范畴，
  并讨论对偶性的作用。（首先出现在 [4] 中。）
\item 在第 7 章，我们对带有哈密顿（Hamiltonian）$U(1)$ 作用的一般 Kähler 靶空间 $X$ 构造一个
  修正 A-模型 TFT，并探讨其性质。
\item 在第 8 章，我们把 Vafa-Witten 理论在亏格 $g_C\ge 2$ 的黎曼面 $C$ 上紧化为
  以 $\cM_H(G,C)$ 为靶空间的修正 A-模型。（Vafa-Witten 理论在黎曼面上的紧化曾在 [20]
  的一个附录中被简要论及；在第 8 章我们完整详细地做这一紧化。）
\item 在第 9 章，利用广义几何与中性号差几何的概念，我们对各种 2d $\sigma$-模型的
  拓扑缺陷进行分类：玻色 $\sigma$-模型，以及 (0,1)、(0,2) 与 (2,2) 超对称模型。
  （首先出现在 [5] 中。）
\end{itemize}

% [Tutor Note: p.7 “凝聚层的等变导出范畴”对应原文 "the equivariant-derived category of coherent
%  sheaves"，见附录 A 的 §A.3；其要点是：带规范场的 B-模型在 $G$ 作用下的膜范畴由
%  $[X/G]$ 上的等变凝聚层的（导出）范畴给出，这是 §A.3 中构造的规范场化作用于
%  $\cD^b(\mathrm{Coh})$ 的结果。]

此外，我们还整理了以下背景材料：

% [原文 p.8]

\begin{itemize}[leftmargin=1.8em,itemsep=3pt,topsep=3pt]
\item 第 2 章回顾 $\cN=4$ SYM 及其各种拓扑扭转版本的对偶性的相关方面。
\item 第 4 章简要讲授范畴、2-范畴以及它们与 TFT 的关系。
\item 附录 C 回顾相伴纤维丛的数学，并解释它们与等变上同调的联系。
\end{itemize}

% [Tutor Note: p.8 此处“tutorial”指 §4.1–§4.3 中作者对范畴语言的自洽讲解（对象/态射/函子、
%  2-范畴、以及边界条件如何构成范畴），它不是对已有文献的直接摘录，而是为后文 §4.4 与第 6 章
%  的 2-范畴陈述铺路；读这些内容时建议与附录 C 的等变上同调回顾合看。]
